\documentclass[10pt,a4paper]{amsart}
\usepackage{amsmath,amssymb,amsthm,graphicx}
\usepackage{verbatim}
\newenvironment{reference}{\comment}{\endcomment}
%\newenvironment{reference}{}{}
\newenvironment{slogan}{\comment}{\endcomment}
\newenvironment{history}{\comment}{\endcomment}

% For commutative diagrams we use Xy-pic
\usepackage[all]{xy}

% We use 2cell for 2-commutative diagrams.
\xyoption{2cell}
% Theorem environments.
%
\theoremstyle{plain}
\newtheorem{theorem}[subsection]{Theorem}
\newtheorem{proposition}[subsection]{Proposition}
\newtheorem{lemma}[subsection]{Lemma}

\theoremstyle{definition}
\newtheorem{definition}[subsection]{Definition}
\newtheorem{example}[subsection]{Example}
\newtheorem{exercise}[subsection]{Exercise}
\newtheorem{situation}[subsection]{Situation}

\theoremstyle{remark}
\newtheorem{remark}[subsection]{Remark}
\newtheorem{remarks}[subsection]{Remarks}

\numberwithin{equation}{subsection}

% Macros
%
\def\lim{\mathop{\mathrm{lim}}\nolimits}
\def\colim{\mathop{\mathrm{colim}}\nolimits}
\def\Spec{\mathop{\mathrm{Spec}}}
\def\Hom{\mathop{\mathrm{Hom}}\nolimits}
\def\Ext{\mathop{\mathrm{Ext}}\nolimits}
\def\SheafHom{\mathop{\mathcal{H}\!\mathit{om}}\nolimits}
\def\SheafExt{\mathop{\mathcal{E}\!\mathit{xt}}\nolimits}
\def\Sch{\mathit{Sch}}
\def\Mor{\mathop{\mathrm{Mor}}\nolimits}
\def\Ob{\mathop{\mathrm{Ob}}\nolimits}
\def\Sh{\mathop{\mathit{Sh}}\nolimits}
\def\NL{\mathop{N\!L}\nolimits}
\def\CH{\mathop{\mathrm{CH}}\nolimits}
\def\proetale{{pro\text{-}\acute{e}tale}}
\def\etale{{\acute{e}tale}}
\def\QCoh{\mathit{QCoh}}
\def\Ker{\mathop{\mathrm{Ker}}}
\def\Im{\mathop{\mathrm{Im}}}
\def\Coker{\mathop{\mathrm{Coker}}}
\def\Coim{\mathop{\mathrm{Coim}}}

% Boxtimes
%
\DeclareMathSymbol{\boxtimes}{\mathbin}{AMSa}{"02}

%
% Macros for moduli stacks/spaces
%
\def\QCohstack{\mathcal{QC}\!\mathit{oh}}
\def\Cohstack{\mathcal{C}\!\mathit{oh}}
\def\Spacesstack{\mathcal{S}\!\mathit{paces}}
\def\Quotfunctor{\mathrm{Quot}}
\def\Hilbfunctor{\mathrm{Hilb}}
\def\Curvesstack{\mathcal{C}\!\mathit{urves}}
\def\Polarizedstack{\mathcal{P}\!\mathit{olarized}}
\def\Complexesstack{\mathcal{C}\!\mathit{omplexes}}
% \Pic is the operator that assigns to X its picard group, usage \Pic(X)
% \Picardstack_{X/B} denotes the Picard stack of X over B
% \Picardfunctor_{X/B} denotes the Picard functor of X over B
\def\Pic{\mathop{\mathrm{Pic}}\nolimits}
\def\Picardstack{\mathcal{P}\!\mathit{ic}}
\def\Picardfunctor{\mathrm{Pic}}
\def\Deformationcategory{\mathcal{D}\!\mathit{ef}}
\usepackage{cleveref}
\crefname{originallemma}{original lemma}{original lemmas}
\Crefname{originallemma}{Original lemma}{Original lemmas}
\makeatletter
\newlabel{divisors-lemma-S1-no-embedded}{{0346}{source}{}{sourcepointer-1}{}}
\newlabel{divisors-lemma-S1-no-embedded@cref}{{[originallemma][0][]0346}{[1][0][]source}{}{}{}}
\makeatother
\crefname{originallemma}{original lemma}{original lemmas}
\Crefname{originallemma}{Original lemma}{Original lemmas}
\makeatletter
\newlabel{divisors-lemma-remove-embedded-points}{{02OL}{source}{}{sourcepointer-2}{}}
\newlabel{divisors-lemma-remove-embedded-points@cref}{{[originallemma][0][]02OL}{[1][0][]source}{}{}{}}
\makeatother
\crefname{originallemma}{original lemma}{original lemmas}
\Crefname{originallemma}{Original lemma}{Original lemmas}
\makeatletter
\newlabel{varieties-lemma-chi-tensor-finite}{{0AYT}{source}{}{sourcepointer-3}{}}
\newlabel{varieties-lemma-chi-tensor-finite@cref}{{[originallemma][0][]0AYT}{[1][0][]source}{}{}{}}
\makeatother
\title{A proper curve with H\textasciicircum{}0(O\_X) = k is connected, Cohen-Macaulay and equidimensional}
\author{The Stacks Project Authors}
\date{stacks-project:curves.tex:local-snapshot-20261012; GNU FDL 1.2 or later}
\begin{document}
\maketitle
\noindent\textit{Editorial note.} All delivered mathematical text is the author's own, copied from the pinned Stacks Project chapter file sources/stacks-curves/pinned-original-source.tex; no statement or proof step was rewritten, reordered or removed.\par\medskip
\noindent\textit{Editorial note.} The counted proof is the author's complete argument: \par\medskip
\noindent\textit{Editorial note.} The label lemma-automatic names the counted statement, and the proof environment immediately following it in the source is delivered verbatim as its proof.\par\medskip
\noindent\textit{Editorial note.} The counted statement is carried by the author's own lemma environment in the chapter source.\par\medskip
\noindent\textit{Allowed standard prerequisites (editorial).}\begin{itemize}
\item The maximal submodule supported in closed points. Divisors, Lemma \textbackslash{}label\{divisors-lemma-remove-embedded-points\}. Proof lines 1289-1297: the maximal coherent submodule I of O\_X supported in closed points exists, and its vanishing is equivalent to X having no embedded points.
\item Sheaves with zero-dimensional support are globally generated. Varieties, Lemma \textbackslash{}label\{varieties-lemma-chi-tensor-finite\}. Proof lines 1291-1293: I is globally generated; since 1 in Gamma(X, O\_X) is not a section of I, one gets I = 0.
\item No embedded points gives Serre's S1. Divisors, Lemma \textbackslash{}label\{divisors-lemma-S1-no-embedded\}. Proof lines 1297-1299: X has (S\_1) because it has no embedded points, hence is Cohen-Macaulay in dimension 1.
\end{itemize}\medskip
\ifdefined\hypertarget\hypertarget{sourcepointer-1}{}\fi
\noindent\textit{Source reference.} Original lemma 0346: sources/stacks-chow/tags.txt, source lines 4081--4081. Prerequisite pointer invoked at proof line(s) 1299; it names a standard fact of a sibling chapter that the counted argument uses, and every step of the counted proof built on it is retained verbatim in the delivered body.\par\medskip
\ifdefined\hypertarget\hypertarget{sourcepointer-2}{}\fi
\noindent\textit{Source reference.} Original lemma 02OL: sources/stacks-chow/tags.txt, source lines 3520--3520. Prerequisite pointer invoked at proof line(s) 1291; it names a standard fact of a sibling chapter that the counted argument uses, and every step of the counted proof built on it is retained verbatim in the delivered body.\par\medskip
\ifdefined\hypertarget\hypertarget{sourcepointer-3}{}\fi
\noindent\textit{Source reference.} Original lemma 0AYT: sources/stacks-chow/tags.txt, source lines 13801--13801. Prerequisite pointer invoked at proof line(s) 1293; it names a standard fact of a sibling chapter that the counted argument uses, and every step of the counted proof built on it is retained verbatim in the delivered body.\par\medskip
\begin{lemma}
\label{lemma-automatic}
Let $k$ be a field. Let $X$ be a proper scheme over $k$ having dimension $1$
and $H^0(X, \mathcal{O}_X) = k$. Then $X$ is connected, Cohen-Macaulay,
and equidimensional of dimension $1$.
\end{lemma}

\begin{proof}
Since $\Gamma(X, \mathcal{O}_X) = k$ has no nontrivial idempotents,
we see that $X$ is connected. This already shows that $X$ is
equidimensional of dimension $1$ (any irreducible component
of dimension $0$ would be a connected component).
Let $\mathcal{I} \subset \mathcal{O}_X$
be the maximal coherent submodule supported in closed points.
Then $\mathcal{I}$ exists
(Divisors, Lemma \ref{divisors-lemma-remove-embedded-points})
and is globally generated
(Varieties, Lemma \ref{varieties-lemma-chi-tensor-finite}).
Since $1 \in \Gamma(X, \mathcal{O}_X)$ is not a section
of $\mathcal{I}$ we conclude that $\mathcal{I} = 0$.
Thus $X$ does not have embedded points
(Divisors, Lemma \ref{divisors-lemma-remove-embedded-points}).
Thus $X$ has $(S_1)$ by
Divisors, Lemma \ref{divisors-lemma-S1-no-embedded}.
Hence $X$ is Cohen-Macaulay.
\end{proof}

\end{document}
